\subsubsection{Modular inverse}
\begin{lstlisting}[style=mystyle, language=C++]
// TC: O(log mod)
// usa: 1/a mod p (con p primo), o por extended gcd si no
#include <bits/stdc++.h>
using namespace std;

// when MOD is prime
long long modInv(long long a, long long mod){
	return modpow(a, mod - 2, mod);
}

// general version using extended gcd (works for any mod coprime with a)
long long extGcd(long long a, long long b, long long& x, long long& y){
	if(b == 0){ x = 1; y = 0; return a; }
	long long x1, y1;
	long long g = extGcd(b, a % b, x1, y1);
	x = y1;
	y = x1 - y1 * (a / b);
	return g;
}

long long modInvGeneral(long long a, long long mod){
	long long x, y;
	long long g = extGcd(a, mod, x, y);
	if(g != 1) return -1; // no inverse
	return (x % mod + mod) % mod;
}

int main(){
	// inverse of 3 mod 1e9+7
	cout << modInv(3, 1000000007LL) << "\n";
	// general (via extgcd)
	cout << modInvGeneral(7, 26LL) << "\n";  // 7*15 = 105 = 1 mod 26
	return 0;
}
\end{lstlisting}
